\documentclass[10pt,a4paper]{amsart}
\usepackage[margin=19mm]{geometry}
\usepackage{amsmath,amssymb,mathtools}
\usepackage[scr=rsfs,cal=euler]{mathalfa}
\usepackage{xfrac}
\usepackage[unicode,hidelinks,bookmarks=false]{hyperref}
\newcommand{\ie}{i.e.\ }
\newcommand{\cf}{cf.\ }
\newcommand{\resp}{respectively\ }
\newcommand{\Subsection}{\subsection}
\providecommand{\nobreakdash}{\nobreak\hbox{-}}
\setlength{\emergencystretch}{2em}
\multlinegap0pt
\usepackage{amssymb}
\usepackage{mathtools}
\usepackage{braket}
\usepackage{dsfont}
\usepackage{enumerate}
\usepackage{tikz}
\newtheorem{theorem}{Theorem}
\newcommand{\IOm}{I\times \Om}
\newcommand{\IOprod}[1]{\left( #1 \right)_{I\times \Omega}}
\newcommand{\ImOm}{I_m\times \Om}
\newcommand{\Om}{\Omega}
\newcommand{\Oprod}[1]{\left(#1\right)_{\Omega}}
\newcommand{\na}{\nabla}
\newcommand{\revB}[1]{\textcolor{black}{#1}}
\newcommand{\vertiii}[1]{{\left\vert\kern-0.25ex\left\vert\kern-0.25ex\left\vert #1
    \right\vert\kern-0.25ex\right\vert\kern-0.25ex\right\vert}}
\title{Fully discrete error analysis of finite element discretizations of time-dependent Stokes equations in a stream-function formulation}
\author{Dmitriy Leykekhman, Boris Vexler, Jakob Wagner}
\date{Source: arXiv:2508.06235v4 (CC BY 4.0)}
\begin{document}
\maketitle

\noindent\emph{Source and licence.} Fully discrete error analysis of finite element discretizations of time-dependent Stokes equations in a stream-function formulation. Version arXiv:2508.06235v4 (CC BY 4.0), distributed by arXiv under the Creative Commons Attribution 4.0 International licence, http://creativecommons.org/licenses/by/4.0/, as recorded in the arXiv licence metadata for this exact version and shown on the abstract page https://arxiv.org/abs/2508.06235v4. Full licence notice: \texttt{licenses/arxiv-2508.06235-CC-BY-4.0.txt}.\par\smallskip

\noindent\textit{Editorial scope.} Counted unit: the article's theorem fully (source0.tex lines 1502--1514). The article's complete original proof of it is delivered with it (source0.tex lines 1515--1650, one \texttt{proof} environment of 136 lines). No mathematics is written, repaired or re-derived: the statement and the proof are the article's own lines, in the article's own notation, and no retained line is substituted or re-flowed.  Retained uncounted as the source-local context the proof needs: lines 1480--1482; lines 1486--1488.  Nothing was omitted from any retained span, so the package carries no omission record.  No retained line contains a citation, so the wrapper carries no bibliography and no bibliography entry.  No retained line contains a figure environment, an image-inclusion command or drawing code, so no image is delivered and the package declares no asset dependency.  Editorial formatting only, in addition: an AMS \texttt{amsart} preamble that restates the article's own declarations and macros used by the retained lines, namely the environments \texttt{theorem} on separate counters, together with the standard packages those lines need.  The retained environments print in this document as Theorem 1 in order of appearance, whereas the article prints the same items with the numbering of its own full document.  The complete native source of the article as distributed in the arXiv e-print for this exact version is bundled unchanged as \texttt{sources/182921/source0.tex}.
\begin{equation}\label{eq:jump_identity ah}
a_h([w]_{m-1}, w_{m-1}^+)=\frac{1}{2}a_h(w^+_{m-1}, w_{m-1}^+)+\frac{1}{2}a_h([w]_{m-1}, [w]_{m-1})-\frac{1}{2}a_h(w_{m-1}^-, w_{m-1}^-),
\end{equation}

\begin{equation}\label{eq:jump_at_0_convention}
    \psi_{kh,0}^- := \Pi_h \psi_0, \text{ and correspondingly } [\psi_{kh}]_0 = \psi_{kh,0}^+ - \Pi_h \psi_0,
\end{equation}

\begin{theorem}\label{lem: stability fully}
  Let $g \in L^2(I;H^1(\Om)) ,\psi_0 \in H^1_0(\Om)$, and let $\psi_{kh} \in X^r_k(V_h)$ be the 
    fully discrete solution to
  \begin{equation}
      {B}_h(\psi_{kh},\varphi_{kh})=\IOprod{\nabla g, \nabla \varphi_{kh}}
      +\Oprod{\na \psi_0,\na \varphi_{kh,0}^+}
      \quad \text{for all $\varphi_{kh}\in X_{k}^r(V_h)$}.
  \end{equation}
  Then there exists a constant $C$ independent of $k$ and $h$ such that 
$$
\sum_{m=1}^M\|\na\partial_t \psi_{kh}\|^2_{I_m\times \Om}+\|\na A_h \psi_{kh}\|^2_{I\times \Om}+\sum_{m=1}^M k^{-1}_m\|[\na \psi_{kh}]_{m-1}\|^2_\Omega + \revB{\vertiii{\psi_{kh,M}^-}^2_h}\le C (\vertiii{\Pi_h \psi_0}^2_h + \|\na g\|^2_{I\times \Om}).
$$  
\end{theorem}

\begin{proof}
  First we test the equation for each $m=1,2,...,M$ with $\varphi_{kh}|_{I_m} =A_h\psi_{kh}|_{I_m}$ and set 
  $\varphi_{kh}$ to zero on all other subintervals. Then on each time subinterval $I_m$, we obtain
$$
\int_{I_m}(\na \partial_t \psi_{kh},\na A_h\psi_{kh} )_\Omega dt+\int_{I_m}\|\na A_h \psi_{kh}\|^2dt +(\na [\psi_{kh}]_{m-1}, \na A_h \psi_{kh,m-1}^+)_\Om=\int_{I_m}(\na g,\na A_h \psi_{kh})_\Omega dt,
$$
where for $m=1$, we have made use of our convention \eqref{eq:jump_at_0_convention}.
% Using that
We may rewrite this, using
$$
(\partial_t \na \psi_{kh},\na A_h \psi_{kh})_{\ImOm} =\int_{I_m}\frac{1}{2}\frac{d}{dt}a_h(\psi_{kh},\psi_{kh})\ dt = \frac{1}{2}a_h(\psi_{kh,m}^-,\psi_{kh,m}^-)-\frac{1}{2}a_h(\psi_{kh,m-1}^+,\psi_{kh,m-1}^+)
$$
and
%since 
$$
(\na [\psi_{kh}]_{m-1}, \na A_h \psi_{kh,m-1}^+)_\Om=a_h([\psi_{kh}]_{m-1},\psi^+_{kh,m-1}),
$$
together with the identity \eqref{eq:jump_identity ah}.
%applying the identity \eqref{eq:jump_identity ah},
By noticing that the term with $a_h(\psi_{kh,m-1}^+,\psi_{kh,m-1}^+)$ cancels out, as a result
we obtain
$$
\begin{aligned}
\frac{1}{2}a_h(\psi_{kh,m}^-,\psi_{kh,m}^-)-\frac{1}{2}a_h(\psi_{kh,m-1}^-,\psi_{kh,m-1}^-)+\frac{1}{2}a_h([\psi_{kh}]_{m-1},[\psi_{kh}]_{m-1}) &+\int_{I_m}\|\na A_h \psi_{kh}\|_\Om^2dt\\&=(\na g, \na A_h \psi_{kh})_{\ImOm}.
\end{aligned}
$$
Summing over $m$, and using the convention $\psi_{kh,0}^- = \Pi_h \psi_0$ yields
$$
\begin{aligned}
\frac{1}{2}a_h(\psi_{kh,M}^-,\psi_{kh,M}^-)+\frac{1}{2}\sum_{m=1}^Ma_h([\psi_{kh}]_{m-1},[\psi_{kh}]_{m-1})&+\int_I \|\na A_h \psi_{kh}\|_\Om^2dt\\
&=(\na g,\na A_h\psi_{kh})_{\IOm}+\frac{1}{2}a_h(\Pi_h\psi_0,\Pi_h\psi_0).
\end{aligned}
$$
Using the Cauchy-Schwarz and Young's inequalities, we have
\revB{
$$
\begin{aligned}
\frac{1}{2} \vertiii{\psi_{kh,M}^-}^2_h + \|\na A_h \psi_{kh}\|^2_{\IOm} &\le \|\na g\|_{I\times \Om}\|\na A_h\psi_{kh}\|_{\IOm}+\frac{1}{2}\vertiii{\Pi_h\psi_0}_h^2\\
&\le \frac{1}{2}\|\na g\|_{I\times \Om}^2+\frac{1}{2}\|\na A_h\psi_{kh}\|^2_{\IOm}+\frac{1}{2}\vertiii{\Pi_h\psi_0}_h^2.
\end{aligned}
 $$
}
Canceling, we obtain
\begin{equation}\label{eq: estimates for integral of ah}
\revB{\vertiii{\psi_{kh,M}^-}^2_h}+ \|\na A_h \psi_{kh}\|^2_{\IOm}\le \|\na g\|^2_{I\times \Om}+\vertiii{\Pi_h\psi_0}_h^2.
\end{equation}


{\it Step 2. Time derivative terms.}

Note that for the case $r=0$ there is nothing to prove here, hence we assume in this step $r\ge 1$.
Noticing that $(t-t_{m-1})\partial_t \psi_{kh}\in \mathbb{P}_r(I_m;V_h)$, we  choose $\varphi_{kh}=(t-t_{m-1})\partial_t \psi_{kh}$ on the time interval $I_m$ and zero on others. Noticing that the term with jumps vanishes, we have
$$
\int_{I_m}(t-t_{m-1})\|\na \partial_t\psi_{kh}\|^2_\Omega dt+\int_{I_m}(t-t_{m-1})a_h(\psi_{kh},\partial_t \psi_{kh})dt =\int_{I_m}(t-t_{m-1})(\na g,\na \partial_t\psi_{kh})_\Omega dt.
$$
For the second term, using the Cauchy-Schwarz and Young's inequalities,
$$
\begin{aligned}
\int_{I_m}(t-t_{m-1})a_h(\psi_{kh},\partial_t \psi_{kh})dt&=\int_{I_m}(t-t_{m-1})( \na \partial_t\psi_{kh},\na A_h \psi_{kh})_\Omega dt\\
&\le \left(\int_{I_m}(t-t_{m-1})\| \na\partial_t \psi_{kh}\|_\Om^2dt\right)^{1/2}\left(\int_{I_m}(t-t_{m-1})\| \na A_h\psi_{kh}\|_\Om^2dt\right)^{1/2}\\
&\le \frac{1}{4}\int_{I_m}(t-t_{m-1})\| \na \partial_t \psi_{kh}\|_\Om^2dt+\int_{I_m}(t-t_{m-1})\| \na A_h\psi_{kh}\|_\Om^2dt.
\end{aligned}
$$
Similarly, 
using the Cauchy-Schwarz and Young's inequalities,
$$
\begin{aligned}
\int_{I_m}(t-t_{m-1})(\na g,\na\partial_t \psi_{kh})_\Omega dt&\le \left(\int_{I_m}(t-t_{m-1})\|\na g\|^2_\Omega dt\right)^{1/2}\left(\int_{I_m}(t-t_{m-1})\|\na \partial_t\psi_{kh}\|^2_\Omega dt\right)^{1/2}\\
&\le \frac{1}{4}\int_{I_m}(t-t_{m-1})\| \na\partial_t \psi_{kh}\|_\Om^2dt+\int_{I_m}(t-t_{m-1})\| \na g\|^2_\Om dt.
\end{aligned}
$$
Combining estimates, and kicking back, we obtain
$$
\int_{I_m}(t-t_{m-1})\|\na \partial_t\psi_{kh}\|^2_\Omega dt\le C\int_{I_m}(t-t_{m-1})\|\na g\|^2_\Omega dt+C\int_{I_m}(t-t_{m-1})\| \na A_h\psi_{kh}\|^2_\Om dt.
$$
Using the inverse inequality
$$
 \int_{I_m}\|\chi_{k}\|^2_\Omega dt\le Ck_{m}^{-1}\int_{I_m}(t-t_{m-1})\|\chi_{k}\|^2_\Omega dt
$$
which holds for any piecewise polynomial in time with fixed maximal degree,
for example for $\chi_k=\na\partial_t \psi_{kh}$, we have
$$
\begin{aligned}
\int_{I_m}\|\na \partial_t\psi_{kh}\|^2_\Omega dt&\le  Ck_{m}^{-1}\int_{I_m}(t-t_{m-1})\|\na g\|^2_\Omega dt+Ck_{m}^{-1}\int_{I_m}(t-t_{m-1})\| \na A_h\psi_{kh}\|_\Om^2dt\\
&\le  C\int_{I_m}\|\na g\|^2_\Omega dt+C\int_{I_m}\| \na A_h\psi_{kh}\|_\Om^2dt.
\end{aligned}
$$
Summing over $m$ and using \eqref{eq: estimates for integral of ah} yields
\begin{equation}\label{eq: estimates for integral of pa_tvkh}
  \int_{I}\|\na \partial_t\psi_{kh}\|^2_\Omega dt\le C(\vertiii{\Pi_h \psi_0}^2_h + \|\na g\|^2_{I\times \Om}).
\end{equation}

{\it Step 3. Jump terms.}


For every $m=1,...,M$, we choose $\varphi_{kh}|_{I_m} =[\psi_{kh}]_{m-1}$ and set $\varphi_{kh}$
to zero on all other subintervals. Inserting this choice yields
$$
\int_{I_m}(\na\partial_t\psi_{kh},[\na \psi_{kh}]_{m-1})_\Om dt+\int_{I_m}(\na A_h\psi_{kh},[\na \psi_{kh}]_{m-1})_\Om dt +\|[\na \psi_{kh}]_{m-1}\|^2_\Omega =\int_{I_m} (\na g, [\na \psi_{kh}]_{m-1})_\Om dt.
$$
Using the Cauchy-Schwarz and Young's inequalities,
$$
\begin{aligned}
\int_{I_m}(\na\partial_t\psi_{kh},[\na \psi_{kh}]_{m-1})_\Om dt&\le 
\left(k_m\int_{I_m}\|\na\partial_t\psi_{kh}\|^2_\Omega dt\right)^{1/2}\left(k^{-1}_m\int_{I_m}\|[\na \psi_{kh}]_{m-1} \|^2_\Omega dt\right)^{1/2}\\
&\le 
k_m\int_{I_m}\|\na\partial_t\psi_{kh}\|^2_\Omega dt+\frac{1}{4}\|[\na \psi_{kh}]_{m-1} \|^2_\Omega,
\end{aligned}
$$
where we used that $\varphi_{kh}|_{I_m} = [\psi_{kh}]_{m-1}$ is constant on $I_m$. Similarly,
$$
\begin{aligned}
\int_{I_m}(\na A_h\psi_{kh},[\na \psi_{kh}]_{m-1})_\Om dt&\le 
\left(k_m\int_{I_m}\|\na A_h\psi_{kh}\|^2_\Omega dt\right)^{1/2}\left(k^{-1}_m\int_{I_m}\|[\na \psi_{kh}]_{m-1} \|^2_\Omega dt\right)^{1/2}\\
&\le 
k_m\int_{I_m}\|\na A_h\psi_{kh}\|^2_\Omega dt+\frac{1}{4}\|[\na \psi_{kh}]_{m-1} \|^2_\Omega,
\end{aligned}
$$
and
$$
\begin{aligned}
  \int_{I_m} (\na g, [\na \psi_{kh}]_{m-1})_\Om dt&\le \left(k_m\int_{I_m}\|\na g\|^2_\Omega dt\right)^{1/2}\left(k^{-1}_m\int_{I_m}\|[\na \psi_{kh}]_{m-1} \|^2_\Omega dt\right)^{1/2}\\
                                             &\le k_m\int_{I_m}\|\na g\|^2_\Omega dt+\frac{1}{4}\|[\na \psi_{kh}]_{m-1} \|^2_\Omega.
\end{aligned}
$$
Combining, we obtain
$$
k^{-1}_m\|[\na \psi_{kh}]_{m-1}\|^2_\Omega \le C\left(\int_{I_m}\|\na g\|^2_\Omega dt+\int_{I_m}\|\na\partial_t\psi_{kh}\|^2_\Omega dt+\int_{I_m}\|\na A_h\psi_{kh}\|^2_\Omega dt\right).
$$ 
Summing over $m$ and using \eqref{eq: estimates for integral of ah} and \eqref{eq: estimates for integral of pa_tvkh}, we obtain
$$
\sum_{m=1}^Mk^{-1}_m\|[\na \psi_{kh}]_{m-1}\|^2_\Omega 
\le C(\vertiii{\Pi_h \psi_0}^2_h + \|\na g\|^2_{I\times \Om}),
$$
which completes the proof of the lemma. 
\end{proof}

\end{document}
